% ----------------------------------------------------------------------------
%                                  ABSTRACT
% ----------------------------------------------------------------------------

\begin{resumo}[ABSTRACT]

    \begin{SingleSpacing}
    
    The continuous monitoring of patients with Parkinson’s disease requires
efficient computer systems to support clinical decision-making. The Painel Parkinson system,
originally developed as a web application, assists healthcare professionals
through informative dashboards. However, the exclusive reliance on web browsers
and the centralization of processing in the cloud result in performance limitations, latency
in rendering complex data, and high operational hosting costs. In light
of this, this study proposes the architectural migration of the Parkinson Panel to a
desktop application. Using the Electron.js framework in conjunction with VueJS on the front-end and Python
with PostgreSQL on the back-end, the project aims to decentralize data processing and
optimize the operational infrastructure. Methodologically, the investigation is classified as
applied, exploratory, and descriptive research with a qualitative approach, grounded
in literature review to support the development of the technological product
designed for the clinical environment. The new desktop version is expected to provide greater agility
in the interpretation of clinical data, improve the user experience, and enable a reduction in
server costs, ensuring more accessible clinical monitoring that is prepared
for future developments. 

        
        \textbf{Keywords}: Parkinson’s disease; Desktop application; Cost optimization; Applied research; Digital health.
    \end{SingleSpacing}
    
\end{resumo}
